\renewcommand{\Anno}{2009}

\renewcommand{\NumeroProblema}{1}

\renewcommand{\Testo}{
Consideriamo l'equazione
$$x^{2}-y^{2}=2^{k}.$$

Determinare i valori \emph{interi} del parametro $k$ per cui esistono esattamente 2009 coppie $(x,y)$ di \emph{interi positivi} che
soddisfano l'equazione.}

\renewcommand{\Risposta}{
$k=4019$ e $k=4010$
}

\renewcommand{\Aiutino}{
Per prima cosa conviene riscrivere l'equazione nella forma $(x+y)(x-y)=2^{k}$. Ora i due fattori devono essere potenze di 2, diciamo $x-y=2^{a}$ e $x+y=2^{b}$ (con quali limitazioni su $a$ e $b$?). Da qui si ricavano i valori di $x$ e $y$.
}

\renewcommand{\SoluzioneUfficiale}{\TitoloSoluzione 
Gli unici valori di $k$ per cui l'equazione ammette esattamente 2009 coppie di soluzioni $(x,y)$ positive sono $k=4019$ e $k=4020$.

Intanto è facile vedere che l'equazione non ha soluzioni intere positive quando $k<0$ (in tal caso infatti $x^{2}-y^{2}$ è intero, mentre $2^{k}$ non lo è) e quando $k=0$ (non ci sono infatti due quadrati positivi che distano 1).  Possiamo pertanto supporre $k>0$. In tal caso, fattorizzando l'equazione
$$(x+y)(x-y)=2^{k},$$
deduciamo che $x+y$ e $x-y$ devono essere entrambi potenze di 2. Possiamo pertanto scrivere
\begin{equation}
	x-y=2^{a},
	\hspace{3em}
	x+y=2^{k-a}.
	\label{eq:sum-sub}
\end{equation}

Essendo $x$ e $y$ positivi, sarà $x-y<x+y$, da cui $a<k-a$,
cioè $a<k/2$. Risolvendo il sistema (\ref{eq:sum-sub}) si ricava
facilmente che
\begin{equation}
	x=\frac{2^{k-a}+2^{a}}{2},
	\hspace{3em}
	y=\frac{2^{k-a}-2^{a}}{2}.
	\label{eq:xy}
\end{equation}

Il termine $2^{k-a}$ è sicuramente pari in quanto $k-a$ è positivo. Affinché $x$ ed $y$ siano interi serve quindi che $2^{a}$ sia pari, da cui $a\geq 1$. Abbiamo cos\`{\i} dimostrato che, per un certo $k>0$ fissato, le soluzioni sono tutte e sole quelle date
dalla (\ref{eq:xy}) con
$$1\leq a<\frac{k}{2}.$$

Per avere esattamente 2009 soluzioni serve quindi che esistano esattamente 2009 valori di $a$ con questa proprietà, il che accade se e solo se
$$2009<\frac{k}{2}\leq 2010,$$
cioè se e solo se $k=2009\cdot 2+1=4019$ oppure $k=2010\cdot
2=4020$.
}

\renewcommand{\SoluzioniAlternative}{}

\renewcommand{\Osservazioni}{
}

\renewcommand{\Commenti}{
}